\section{模型检验、敏感性与适用范围}
\subsection{分层验证}
问题一的配比模型以 A4--A5 训练、A6--A7 的 1M 记录作独立检验；A12--A15 大尺度估计仅用于压力测试。问题二 M0 按整条参数规模轨迹留出，M1 按完整 $(N,D)$ 单元分组，$D_B=600$ 和 B7 新质量水平单列，避免同一单元同时进入拟合与验证。问题三对各预算回代原始 FLOPs、支持域与活跃约束，比较三种题设质量成本、基线质量、上下文长度及 B6 离散质量档位。问题四按家族分组和时间留出，另核查开放状态、C1/C4 连接、历史评分协议及 C8 叶任务标尺。

\begin{table}[htbp]\centering\caption{关键验证数值与解释边界}
\small\begin{tabular}{p{3.2cm}p{2.8cm}p{5.7cm}}\toprule
验证 & 结果 & 解释\\\midrule
配比 1M 独立测试 & RMSE 0.23195 & 13 域平均 Loss，不能代表大尺度绝对预测\\
B1 留完整轨迹 & RMSE 0.000151 & 附件内部强规则关系，不代表现实泛化误差\\
B6 单元分组 CV & RMSE 0.05951 & 半合成质量响应的结构选择\\
B7 新质量档位 & RMSE 0.04459 & 限于 B7 的 D$\leq$300B\\
后训练未来留出 & 8.0189 / 11.736 & 仅规模 / 加时间的家族均衡 RMSE\\\bottomrule
\end{tabular}\end{table}

\subsection{共同限制与不得外推之处}
质量评分 $Q_A$ 是按 A1 标尺构造的相对指标，$Q_B$ 是 B6 实验水平；二者没有共同样本、共同 Loss 或共同配比的联合标定。第三问固定的 $p^{\rm ref}$ 因此不改变数值目标和成本，表中仅为 $(N,D,Q_B)$ 的条件最优。问题一二次 ILR 的非加性检验提供模型内局部证据，不证明领域内在互补。问题二逐域规模指数有负值，且只有 13 域目标可估，跨尺度统一衰减不成立。问题四排行榜记录、算力元数据、Loss 测试之间存在样本选择和协议差异，无法构建经验证的长期能力前沿。

为避免不确定性过窄，文中将三类量分开：数据重抽样区间描述当前附件内的不稳定性；支持域扫描描述边界选择；成本函数与质量机制替换描述结构情景。三者不能相加为单一“95\% 预测区间”。尤其 B1 的极窄轨迹 Bootstrap、第三问的条件 Bootstrap 与第四问的宽逻辑界，性质完全不同。
